% Follow-up 1, fixed surfaces. 2026-11-22.
% GPT-6 Astra Pro, /root/fixed_surface. Numerical effort is not exposed.
% This is an input fragment. No claim of proof-assistant verification.

\section{Follow-up 1: the conormal map for specified surfaces}
\label{sec:f1-fixed-surface}

We use the degree convention $\deg_\theta Y=\int_Y\theta^2$.
The original DEG\_L S-02 gives, for a smooth surface $i:Y\hookrightarrow A$
in an abelian fourfold, the dual of Bloch's semiregularity map:
\[
 \mu_i:H^1(A,\mathcal O_A)\otimes\bigwedge^3H^0(A,\Omega_A^1)
       \longrightarrow H^1(Y,N_{Y/A}^{\vee}\otimes K_Y).
 \tag{F1S0.1}
\]
The first source has dimension sixteen. All surfaces below are specified
by actual embeddings. The norm-fibre construction in Schoen (1988) is
discussed separately at the end; its identification with a particular
smooth embedded surface is not assumed.

\subsection{HD-F1-S01: the fourfold map and its cokernel}

Choose a nonzero invariant four-form $\eta$ on $A$, and put
$V=\operatorname{Lie}(A)$. Contraction by $\eta$ identifies
$T_A$ with $\Omega_A^3$. Adjunction identifies $\det N_{Y/A}$ with
$K_Y$, and the rank-two determinant pairing identifies
$N_{Y/A}$ with $N_{Y/A}^{\vee}\otimes K_Y$.
Under these identifications (F1S0.1) is the actual map
\[
 \operatorname{ob}_i:H^1(A,T_A)\xrightarrow{i^*}
 H^1(Y,i^*T_A)\longrightarrow H^1(Y,N_{Y/A}).
 \tag{F1S1.1}
\]
Indeed, in a tangent-normal frame $(u_1,u_2,n_1,n_2)$ with
$\eta(u_1,u_2,n_1,n_2)=1$, the map
$\bigwedge^3i^*\Omega_A^1\to N^{\vee}\otimes K_Y$
annihilates the contractions by $u_1,u_2$ and sends contraction by
$n_j$ to the determinant-pairing image of $n_j$. This checks the
sheaf map, not just an equality of vector-space dimensions.

When $i^*:H^1(A,\mathcal O_A)\to H^1(Y,\mathcal O_Y)$ is an
isomorphism, the tangent-normal long exact sequence gives
\[
 \operatorname{coker}\mu_i\simeq
 \ker\left(H^2(Y,T_Y)\longrightarrow V\otimes H^2(Y,\mathcal O_Y)\right).
 \tag{F1S1.2}
\]
This premise holds for the Albanese embedding used in HD-F1-S05.
Serre duality rewrites (F1S1.2) as
\[
 (\operatorname{coker}\mu_i)^\vee\simeq
 \operatorname{coker}\left(
 H^0(A,\Omega_A^1)\otimes H^0(Y,K_Y)
 \longrightarrow H^0(Y,\Omega_Y^1\otimes K_Y)\right).
 \tag{F1S1.3}
\]
The arrow is restriction of one-forms followed by multiplication by a
canonical form. Thus this displayed multiplication map is a precise
remaining calculation for any such Albanese surface.

The identity (F1S1.1) agrees with Bae--Kool--Park,
\emph{Counting surfaces on Calabi--Yau 4-folds I: Foundations},
arXiv:2208.09474v2, Proposition 4.24, Remark 4.25 and Corollary 4.27.
It is not claimed here as new.

\subsection{HD-F1-S02: the actual graph detectors are semiregular}

Use the matched-product sequence of HD-F1-T02. In particular
$A_k=B_k\times B_k$, $\theta_k=\beta_k\boxplus\beta_k$ is principal,
and $\tau_k:B_k\to B_k$ is the polarization-preserving involution
constructed there. Its semilinearity for the specified $K$ action is
\[
 \tau_k\iota(a)=\iota(\bar a)\tau_k\quad(a\in K),
 \tag{F1S2.1}
\]
and is holomorphic for the complex structure of $B_k$.
Here $K=\mathbb Q(\sqrt{-3})$. Set
\[
 i_k:B_k\hookrightarrow A_k,\quad x\longmapsto(x,\tau_kx),
 \qquad Y_k=i_k(B_k).
\]
The graph has $K_{Y_k}=\mathcal O_{Y_k}$ and trivial normal bundle
$N_k=\mathcal O_{Y_k}\otimes(V_k/U_k)$, where
$V_k=\operatorname{Lie}(A_k)$ and $U_k=\operatorname{Lie}(Y_k)$.
Equation (F1S1.1) becomes
\[
 i_k^*\otimes q_k:
 H^1(A_k,\mathcal O_{A_k})\otimes V_k
 \longrightarrow H^1(Y_k,\mathcal O_{Y_k})\otimes(V_k/U_k).
 \tag{F1S2.2}
\]
The first projection of $A_k$ retracts $i_k$, so $i_k^*$ is
surjective. The quotient $q_k$ is surjective. Consequently
\[
 \operatorname{rank}\mu_{i_k}=4,\quad
 \dim\ker\mu_{i_k}=12,\quad \operatorname{coker}\mu_{i_k}=0.
 \tag{F1S2.3}
\]
Bloch's map is therefore injective, with four-dimensional source.

The polarization restricts to $2\beta_k$, whence
\[
 \deg_{\theta_k}Y_k=(2\beta_k)^2=8.
 \tag{F1S2.4}
\]
To check the detector, work on $H^1(-,K)$ and write its two $K$
eigenspaces with subscripts $+$ and $-$. Pullback of the $+$
eigenspace of $A_k$ is
\[
 H^1(B_k,K)_+\oplus H^1(B_k,K)_+
  \longrightarrow H^1(B_k,K),\quad (u,v)\longmapsto u+\tau_k^*v.
 \tag{F1S2.5}
\]
By (F1S2.1), $\tau_k^*$ maps the second summand isomorphically
onto the $-$ eigenspace. Thus (F1S2.5) is an isomorphism, and its
fourth exterior power has nonzero period on $Y_k$.
The exterior decomposition of $H^4(A_k,K)$ shows that the two
extreme determinant summands pair only with one another under
Poincare duality. Hence $e_W[Y_k]\ne0$. This proves
\[
 \delta_W(A_k,\theta_k)\leq 8.
 \tag{F1S2.6}
\]
The exact lattice/core identities $J_W(s_k)=p^{4k}$ and
$(h_k)_*[\Gamma_0]=p^{2k}[Y_k]$ are supplied by HD-F1-T02.
The statements here retain its specified matched product bases.

\subsection{HD-F1-S03: full-class Hodge locus and the obstruction}

We compute the locus of the \emph{whole} class $[Y_k]$, not just its
already Hodge Weil projection. Suppress $k$ and put
$U=H_1(Y,\mathbb Q)$, $V=H_1(A,\mathbb Q)$ and
$I=\iota(\sqrt{-3})$. Equation (F1S2.1) gives
$IY=\Gamma_{-\tau}$, in the sense of rational homology subspaces.
The product Riemann form satisfies
\[
 V=U\oplus IU,\qquad
 I=\begin{pmatrix}0&-3\,\mathrm{id}\cr\mathrm{id}&0\end{pmatrix},
 \qquad \psi=\begin{pmatrix}\psi_U&0\cr0&3\psi_U\end{pmatrix}.
 \tag{F1S3.1}
\]
For example, the mixed pairing vanishes because it is the sum of
$\psi_B(x,z)$ and $\psi_B(\tau x,-\tau z)$; and
$\psi(Iu,Iv)=3\psi(u,v)$ by the Rosati relation.

The class $[Y]$ is a nonzero decomposable rational four-form on
$V$ whose contraction kernel is precisely $U$. It is of type
$(2,2)$ exactly when $U$ is a Hodge subspace: a $(2,2)$ class is
fixed by the Hodge circle, so its contraction kernel is invariant;
conversely a complex two-dimensional quotient $V/U$ has top
exterior volume of type $(2,2)$. Thus its Hodge locus within the
polarized $K$ period domain consists exactly of the complex
structures
\[
 J=\operatorname{diag}(J_U,J_U)
 \quad\text{on }U\oplus IU,
 \tag{F1S3.2}
\]
where $J_U$ is any complex structure positive for $\psi_U$.
The group morphism producing this locus is
\[
 \operatorname{Sp}(U,\psi_U)\longrightarrow
 \operatorname{Cent}_{\operatorname{Sp}(V,\psi)}(I)
       \simeq U(2,2),\qquad g\longmapsto\operatorname{diag}(g,g).
 \tag{F1S3.3}
\]
Its image domain is the genus-two Siegel domain
$\operatorname{Sp}_4(\mathbb R)/U(2)$, of complex dimension three,
inside the four-dimensional Weil domain. Equations (F1S3.1)--(F1S3.3)
prove the comparison by an actual group morphism and an equality
of its period-domain image with the class Hodge locus.

Here is the infinitesimal map with all its ranks. Let
$F=H^{0,-1}(U)$ and identify $\overline F=H^{-1,0}(U)$ using
chosen dual symplectic bases. Set
$f_\pm(u)=u\mp i Iu/\sqrt3$. The polarized $K$ tangent space
is parametrized by $Z\in M_2(\mathbb C)$ through
\[
 \kappa(f_+(u))=\overline{f_-}(Zu),\qquad
 \kappa(f_-(u))=\overline{f_+}(Z^tu).
 \tag{F1S3.4}
\]
Indeed $K$-linearity forces these two blocks, and infinitesimal
isotropy for $\psi$ makes them transposes. Restricting to
$u=(f_+(u)+f_-(u))/2$ and projecting to the normal summand
$I\overline F$ gives
\[
 \operatorname{ob}_Y|_{T\mathcal D_W}(Z)
       =\frac{i}{2\sqrt3}(Z^t-Z).
 \tag{F1S3.5}
\]
The identification of $I\overline F$ with $\overline F$ is by $I$.
The map has rank one and kernel $\operatorname{Sym}_2(\mathbb C)$,
of dimension three. In particular $Z=\left(\begin{smallmatrix}
0&1\\-1&0\end{smallmatrix}\right)$ has nonzero obstruction.

The standard semiregularity compatibility, with the same choice of
sign in the contraction maps, is
\[
 \sigma_Y\bigl(\operatorname{ob}_Y(\kappa)\bigr)
       =\kappa\mathbin{\lrcorner}[Y]
       \quad\text{in }H^3(A,\Omega_A^1).
 \tag{F1S3.6}
\]
See Buchweitz--Flenner, Corollary 4.12 and Proposition 8.2, as
used in DEG\_L S-04, or Bae--Kool--Park, Proposition 4.24.
By HD-F1-S02 the first arrow is injective. Hence the nonzero
direction in (F1S3.5) is an actual direction where the whole class
ceases to be Hodge although $e_W[Y]$ remains Hodge.
Over the three-dimensional locus (F1S3.2), the fixed rational
subspace $U$ gives an actual family of abelian subvarieties, with
the same degree eight. This accounts precisely for the family
consequence of the graph calculation. It does not extend that
graph Hilbert component across the fourth Weil direction.

\subsection{HD-F1-S04: graph equations in biextensions and K-theory}

For a graph as above define
$q_\tau:B\times B\to B$ by $q_\tau(x,y)=y-\tau x$.
Let $L_\beta$ be a symmetric line bundle representing the principal
polarization, and let $\mathcal P_\beta$ be its normalized
Poincare biextension on $B\times B$. The theorem of the square,
with this normalization, gives the cohomological identity
\[
 c_1(q_\tau^*L_\beta)=\beta_1+\beta_2
        -c_1((\tau\times1)^*\mathcal P_\beta),\qquad
 [Y]=q_\tau^*[0]=\tfrac12(q_\tau^*\beta)^2.
 \tag{F1S4.1}
\]
The graph is a fibre of the smooth quotient $q_\tau$, so this is
an equality of its fundamental class with the pullback of a point.
Since its normal bundle is trivial, Grothendieck--Riemann--Roch
gives, in rational cohomology,
\[
 \operatorname{ch}(i_*\mathcal O_Y)=[Y],\qquad
 \operatorname{ch}(\psi^r[i_*\mathcal O_Y])=r^2[Y],\qquad
 [m]^*[Y]=m^4[Y].
 \tag{F1S4.2}
\]
The first follows from $i_*(\operatorname{td}(N)^{-1})$; the
second is the degree-two Adams character formula. These equations
specify the factors in the comparison of the actual K-class,
the Adams operation, and multiplication on the abelian variety.
Translations give $A/Y\to\operatorname{Hilb}(A)$ with tangent
$\operatorname{Lie}(A/Y)\simeq H^0(Y,N)$, both of dimension two.
Together with (F1S3.2), this supplies actual Hilbert families of
the graphs wherever their full class is Hodge.

\subsection{HD-F1-S05: the smooth Schoen family and its Weil class}

We now use the smooth surfaces obtained by Schoen's 2007 smoothing
construction, with the precise cohomological description in
van Geemen, \emph{Weil Classes and Decomposable Abelian Fourfolds},
SIGMA 18 (2022), 097, Lemma 3.3 and Sections 6.1--6.2.
On $J\times J$, for a genus-two Jacobian $J$ with theta curve $C$,
put $S_0=\Delta_J+(C\times C)$. Use van Geemen's classes
$\omega_1,\omega_2,\omega_\sigma$. They satisfy
\[
 [S_0]=\tfrac12(\omega_1+\omega_2-\omega_\sigma)^2
                       +\omega_1\omega_2.
 \tag{F1S5.1}
\]
Writing $\zeta=(-1+\sqrt{-3})/2$, the specified $K$ representation
on coefficient columns of one-forms is
\[
 \zeta\longmapsto\begin{pmatrix}-1&-1\\1&0\end{pmatrix},\qquad
 \beta=2\omega_1-\omega_\sigma+2\omega_2.
 \tag{F1S5.2}
\]
Van Geemen's Section 3.2 uses $g\cdot\omega=(g^t)^*\omega$.
Thus the actual product-coordinate endomorphism corresponding
to the displayed matrix is its transpose; HD-F1-S07 fixes this
convention explicitly and verifies polarization compatibility.
This polarization has type $(1,1,3,3)$ and $\beta^4=216$.
The eigenforms and determinant plane are
\[
 \omega_+=\zeta^2\omega_1+\zeta\omega_\sigma+\omega_2,
 \quad \omega_-=\zeta\omega_1+\zeta^2\omega_\sigma+\omega_2,
 \quad W_K\otimes K=K\omega_+^2\oplus K\omega_-^2.
 \tag{F1S5.3}
\]
Expanding using $\zeta^2+\zeta+1=0$ gives the explicit rational
decomposition
\[
 [S_0]=\frac16\beta^2+w,\qquad
 w=\frac16(\zeta\omega_+^2+\zeta^2\omega_-^2)\ne0.
 \tag{F1S5.4}
\]
The two extreme eigenforms are independent and have nonzero
coefficients, and Galois conjugation fixes $w$. This proves
$w\in W_K\setminus\{0\}$ directly. The $K$ character of
$\beta^2$ is the middle character, so its determinant-plane
projection is zero. Therefore $e_W[S_0]=w$.
The mixed product $\beta^2w$ vanishes by the same character
decomposition, and consequently
\[
 \deg_\beta S=\beta^2[S]=\frac16\beta^4=36
 \tag{F1S5.5}
\]
for every smooth member $S$ of this deformation family.
The flat family transports (F1S5.4), including its nonzero
determinant component, to those members.

For a smooth Schoen member the Albanese embedding is Lagrangian
for a nondegenerate invariant two-form, and
$q=4$, $p_g=5$, $h^{1,1}=12$; see Ciliberto--Mendes Lopes--Roulleau,
\emph{On Schoen surfaces}, and Bae--Kool--Park, Corollary 4.30
and Example 4.32. We give a direct rank computation which
proves semiregularity for smooth members sufficiently close
to the normal-crossing seed.

Choose a compatible symplectic real coframe on $J$ and write
its holomorphic forms as $dz_1,dz_2$. Use $dz_3,dz_4$ for the
second copy. Normalize $\omega_J=(i/2)\sum dz_j\wedge d\bar z_j$
and $\int_J\omega_J^2=2$. Identify $H^{1,1}(J^2)$ with
$M_4(\mathbb C)$ by
$M\mapsto\alpha_M=\sum M_{ab}dz_a\wedge d\bar z_b$.
For two block matrices
$M=\left(\begin{smallmatrix}A&B\\C&D\end{smallmatrix}\right)$
and $M'$ with two-by-two blocks, put $E=A+B+C+D$ and
$E'=A'+B'+C'+D'$. Direct restriction to the diagonal and
to $C\times C$ gives
\[
 \begin{split}
 -\tfrac14\int_{J^2}\alpha_M\alpha_{M'}[S_0]
 ={}&\operatorname{tr}E\operatorname{tr}E'-\operatorname{tr}(EE')\\
 &+\operatorname{tr}A\operatorname{tr}D'
   +\operatorname{tr}D\operatorname{tr}A'
   -\operatorname{tr}(BC')-\operatorname{tr}(CB').
 \end{split}
 \tag{F1S5.6}
\]
For example $\int_C dz_a\wedge d\bar z_b=-2i\delta_{ab}$,
and $\int_J\alpha_X\alpha_{X'}=-4(\operatorname{tr}X
\operatorname{tr}X'-\operatorname{tr}(XX'))$, which proves
the formula term by term.

Put $L(E)=\operatorname{tr}(E)\mathrm{id}-E$. A matrix is
in the radical of (F1S5.6) exactly when
\[
 L(E)+\operatorname{tr}D\,\mathrm{id}=0,\quad
 L(E)+\operatorname{tr}A\,\mathrm{id}=0,\quad
 L(E)-C=0,\quad L(E)-B=0.
\]
Solving these equations gives precisely
\[
 A\in M_2(\mathbb C),\quad
 B=C=-\operatorname{tr}(A)\mathrm{id},\quad
 D=\operatorname{tr}(A)\mathrm{id}-A.
 \tag{F1S5.7}
\]
Its dimension is four. Therefore the Gram form has rank twelve.
For a separately reproducible arithmetic check, in the
lexicographic basis $dz_a\wedge d\bar z_b$, the principal
minor indexed by
$1,2,3,4,5,6,7,8,9,10,11,13$ has determinant $2^{24}$;
the accompanying script computes the full matrix in exact
exterior algebra.

This nonzero minor persists in an open neighborhood in the
smooth abelian family. Smooth Schoen fibres occur arbitrarily
near the seed by its actual smoothing. On such a smooth
fibre, the Gram form factors as restriction to
$H^{1,1}(S)$, followed by its cup-product pairing. Since
$h^{1,1}(S)=12$, the restriction has rank twelve and is
surjective. Choose such a smooth member. Contraction by its
symplectic form identifies the normal sequence with the
cotangent sequence, so (F1S1.1) is identified with
\[
 H^1(A,\Omega_A^1)\xrightarrow{i^*}H^1(S,\Omega_S^1).
 \tag{F1S5.8}
\]
For the chosen member this is surjective and its target has
dimension twelve. Thus the exact answer for that embedding is
\[
 \operatorname{rank}\mu_i=12,\quad\dim\ker\mu_i=4,
 \quad\operatorname{coker}\mu_i=0.
 \tag{F1S5.9}
\]
The semiregularity is established here from (F1S5.6)--(F1S5.7)
and the actual smoothing; it also agrees with the existence
statement recorded in Bae--Kool--Park, Example 4.32.

For this actual member, the relative Hilbert family dominates
an open subset of the four-dimensional Weil family. Here are
the exact deformation steps. The full class
$[S]=\beta^2/6+w$ is flat and of type $(2,2)$ throughout its
marked Weil base. Ran, \emph{Semiregularity, obstructions and
deformations of Hodge classes}, Ann. Scuola Norm. Sup. Pisa
28 (1999), Theorem 0(ii), pp.~809--810, puts every
small-extension obstruction for the pair $(A,S)$, relative
to ambient deformations preserving this full Hodge class,
in $\ker\sigma_S$. That kernel is zero by (F1S5.9).
Thus the relative Hilbert functor is formally smooth at this
point over the smooth Weil base. The relative Hilbert scheme
is locally of finite presentation, so its projection is
smooth there and has open image. Its universal subscheme is
an actual family of smooth embedded detecting surfaces after
shrinking around this point. This proves the local dominance
needed in HD-F1-S07 directly from the conormal calculation.
In comparison, (F1S3.5) computes a nonzero obstruction for
the graph in the fourth Weil direction; its full-class
Hodge locus is the three-dimensional domain (F1S3.2).

\subsection{HD-F1-S06: smooth Schoen pullbacks that fail S-02}

Fix a smooth Schoen member $S\subset A=\operatorname{Alb}(S)$
as in the preceding subsection. For every integer $m\ge2$ form
the actual fibre square
\[
 \begin{array}{ccc}
 S_m=[m]^{-1}S&\xrightarrow{j_m}&A\\
 \downarrow f_m&&\downarrow[m]\\
 S&\xrightarrow{i}&A.
 \end{array}
 \tag{F1S6.1}
\]
The map $f_m$ is etale of degree $m^8$, and $S_m$ is smooth.
It is connected: the Albanese map induces a surjection
$\pi_1(S)\to H_1(S,\mathbb Z)/\mathrm{tors}=\pi_1(A)$,
so its composition with $\pi_1(A)/m\pi_1(A)$ is onto; this is
the monodromy criterion for the pullback of the connected
cover $[m]$. Hence $S_m$ is a smooth integral embedded surface.

Here is a direct dimension obstruction for its conormal map.
For any smooth surface $Y$ in an abelian fourfold,
\[
 c_1(N)=K_Y,\quad c_2(N)=K_Y^2-c_2(Y),\qquad
 \chi(N)=14\chi(\mathcal O_Y)-2K_Y^2.
 \tag{F1S6.2}
\]
The Chern equations follow from $c(T_Y)c(N)=1$. Riemann--Roch
then gives $\chi(N)=(-5K_Y^2+7c_2(Y))/6$, equal to the last
expression by Noether's formula. For $S$, the Hodge numbers
give $\chi(\mathcal O_S)=2$ and $c_2(S)=8$; Noether gives
$K_S^2=16$. Thus $\chi(N_{S/A})=-4$.
Etale pullback identifies $N_{S_m/A}=f_m^*N_{S/A}$ and gives
\[
 \chi(N_{S_m/A})=-4m^8.
 \tag{F1S6.3}
\]
Rank-two adjunction and Serre duality give
$h^2(N_{S_m/A})=h^0(N_{S_m/A})$. The pulled-back symplectic
form identifies $N_{S_m/A}$ with $\Omega_{S_m}^1$; pullback
of one-forms from $S$ is injective, so $h^0(N_{S_m/A})\ge4$.
Consequently
\[
 h^1(N_{S_m/A})=2h^0(N_{S_m/A})+4m^8\ge8+4m^8,
 \qquad
 \dim\operatorname{coker}\mu_{j_m}\ge4m^8-8>0.
 \tag{F1S6.4}
\]
For $m=2$ these bounds are $h^1(N)\ge1032$ and
$\dim\operatorname{coker}\mu\ge1016$. Thus S-02 fails for
every one of these specified smooth embeddings.
Nevertheless the determinant calculation is exact:
\[
 [S_m]=[m]^*[S]=m^4[S],\qquad e_W[S_m]=m^4w\ne0,
 \qquad \deg_\beta S_m=36m^4.
 \tag{F1S6.5}
\]
These surfaces also occur in an actual family: pull back the
smooth Schoen family by the relative multiplication map $[m]$.
The degrees in (F1S6.5) are constant on its smooth base.
Thus the morphism (F1S6.1) simultaneously exhibits failure of
the conormal surjectivity and persistence of actual detecting
surfaces in that family. No converse from detector existence
to semiregularity is used.

\subsection{HD-F1-S07: the degree-nine quotient to the principal component}

There is an explicit polarized $\mathcal O_K$-linear isogeny from
the type $(1,1,3,3)$ component of HD-F1-S05 to the principal
Eisenstein component containing $E_0^4$ with balanced action.
We construct its lattice at a specified point and then transport
that lattice on the common marked period domain.

Take $J=E_0^2$ with product principal polarization, where
$E_0=\mathbb C/\mathcal O_K$. The class $\beta$ on $J^2$
has symmetric coefficient matrix
\[
 B=\begin{pmatrix}2&-1\\-1&2\end{pmatrix}.
\]
In column coordinates on $J^2$ use the actual endomorphism
\[
 T=\begin{pmatrix}-1&1\\-1&0\end{pmatrix}.
 \tag{F1S7.1}
\]
Its pullback on coefficient columns of one-forms is
$T^t=\left(\begin{smallmatrix}-1&-1\\1&0\end{smallmatrix}\right)$,
the representation in (F1S5.2). In particular
$T^2+T+1=0$ and $T^tBT=B$; this checks the action and
polarization in a single fixed convention.

Regroup the four elliptic coordinates into two blocks, each
consisting of the matching coordinate in the two copies of $J$.
On each block define
\[
 R=\begin{pmatrix}1&\zeta\\1&\zeta^2\end{pmatrix},
 \qquad D=\operatorname{diag}(\zeta^2,\zeta).
\]
Complex conjugate transpose is denoted by a dagger. Direct
multiplication, using $\zeta+\zeta^2=-1$, gives
\[
 R^\dagger R=B,\qquad RT=DR,\qquad
 \det R=\zeta^2-\zeta,\qquad N(\det R)=3.
 \tag{F1S7.2}
\]
Thus the actual homomorphism $q_0=R\oplus R:E_0^4\to E_0^4$
has degree nine, satisfies $q_0^*\theta_0=\beta$, and is
$\mathcal O_K$-linear for $T\oplus T$ on its source and
$D\oplus D$ on its target. Reordering the target factors gives
the balanced action of HD-F1-T02, with the two embeddings
possibly interchanged. Its kernel $H$ has order nine. It is
$\mathcal O_K$-stable by $RT=DR$ and isotropic for the
polarization pairing because $\beta$ is the pullback of a
principal polarization. This proves the needed stability,
rather than assuming it for an arbitrary polarization quotient.

Here is the period-domain map. Identify the target rational
vector space with the source by $q_0$. Its integral lattice
pulls back to a superlattice $\Lambda'$ of $\Lambda$, with
$[\Lambda':\Lambda]=9$, stable under $T\oplus T$. The same
alternating form $\beta$ is unimodular on $\Lambda'$. Every
positive complex structure in the connected polarized Weil
domain for $(\Lambda,\beta,T)$ is also a positive complex
structure for $(\Lambda',\beta,T)$. The maps
\[
 (V_{\mathbb R},J)/\Lambda\longrightarrow
 (V_{\mathbb R},J)/\Lambda'
 \tag{F1S7.3}
\]
are degree-nine polarized $\mathcal O_K$-linear isogenies.
On the marked domains (F1S7.3) is the identity map of a common
four-dimensional domain. On moduli one first passes to the
finite level cover preserving $\Lambda'$; the resulting map is
dominant and generically finite onto the principal component
containing the displayed target $(E_0^4,\theta_0,D\oplus D)$.
This identifies the target component by an actual integral
example and a period-domain morphism.

Push a smooth Schoen surface $S$ through (F1S7.3), and write
$q_*[S]=d[Z]$ with $Z=q(S)$ reduced and integral. The projection
formula gives
\[
 d\deg_\theta Z=\int_Sq^*\theta^2=36.
 \tag{F1S7.4}
\]
The rational $K$-linear cohomology isomorphism $q^*$ identifies
the determinant planes, and $q_*q^*=9\,\mathrm{id}$.
Consequently $e_Wq_*[S]\ne0$, so $e_W[Z]\ne0$ and
$\deg_\theta Z\le36$.
The relative Hilbert family just proved smooth in HD-F1-S05
dominates an open subset of the four-dimensional source
Weil family. The quotient therefore gives these detectors on
a dense subset of the specified principal component. Applying
the proper relative Chow specialization argument already
proved in the intake extends this degree bound over that
component: a specialized effective cycle has total degree at
most thirty-six and a nonzero flat determinant component,
so at least one integral component detects and has degree at
most thirty-six. Thus
\[
 \delta_W(s)\le36
 \quad\text{on this principal Eisenstein component.}
 \tag{F1S7.5}
\]
The descent uses the actual quotient (F1S7.3); no smoothness
assertion about its reduced image $Z$ is needed.

For the original smooth $S$ the reduced image has degree
exactly thirty-six. To see this, let
$H_S=\{h\in\ker q:S+h=S\}$. For an integral $S$, an element
$h$ outside $H_S$ meets a general translate point of $S$
in no second point on $S$: the proper intersection
$S\cap(S+h)$ is a proper closed subset of $S$.
Removing these finitely many subsets shows that the generic
degree of $q|_S$ is $d=|H_S|$. The translation action of
$H_S$ is free, so $S\to S/H_S$ is etale and
$|H_S|$ divides $\chi(\mathcal O_S)=2$. It also divides
$|\ker q|=9$. Hence $d=1$, and (F1S7.4) gives
$\deg_\theta Z=36$. The bound after specialization remains
the inequality (F1S7.5), since special effective cycles can
have several components.

The whole class, which is useful for the division calculation,
is preserved as well. The explicit exterior calculation in
HD-F1-S05 gives $w^2=2$. Put $w'=q_*w$. Then
\[
 [Z]=\frac32\theta^2+w',\qquad
 w'\in W_K\cap H^4(A/\ker q,\mathbb Z),\qquad
 (w')^2=18.
 \tag{F1S7.6}
\]
Indeed $q_*\beta^2=9\theta^2$, and $q^*q_*w=9w$ gives
$\int(q_*w)^2=9\int w^2$. Finally $\theta^2/2$ is integral:
in a symplectic integral basis for the principal polarization,
$\theta=\sum_{j=1}^4x_j\wedge y_j$, so
$\theta^2/2=\sum_{j<l}x_j\wedge y_j\wedge x_l\wedge y_l$.
Thus (F1S7.6) has an integral $w'$. Proper specialization of
the actual relative cycle preserves its entire flat class.
Accordingly every point of the stated principal component
has an effective cycle of degree thirty-six with the form
(F1S7.6), after the appropriate marking/finite cover; the
special cycle may be a sum of integral surfaces.

\subsection{HD-F1-S08: an actual Prym degeneration to the same model}

The polarization induced from the Jacobian on the Prym of an
etale cyclic triple cover of a genus-three curve has type
$(1,1,3,3)$. A direct published locator is Lange--Ortega,
\emph{Prym varieties of cyclic coverings}, arXiv:0805.1020,
Section 4, p.~6: its formula for etale degree $n$ covers is
$D=(1^{p-(g-1)},n^{g-1})$, where $p=(n-1)(g-1)$.
Here is, in addition, an actual degeneration identifying the
polarized $K$ model with (F1S5.2); equality of polarization
types alone is not being used for that identification.

Take a smooth genus-two curve $D$, a marked point $d\in D$,
an elliptic curve $E'$, and a nonzero point $t\in E'[3]$.
Let $e_j=e_0+jt$, $j\in\mathbb Z/3$. Attach the marked point
on a copy $D_j$ of $D$ to $e_j$ on $E'$. The resulting nodal
curve $\widetilde C_0$ has a tree as its dual graph and genus
seven. Translation by $t$ on $E'$ and permutation
$D_j\to D_{j+1}$ define a free cyclic action. The quotient
$C_0$ is the compact-type curve $D\cup E$, where
$E=E'/\langle t\rangle$; it has genus three.

Choose matching local coordinates at the three nodes by the
cyclic action and plumb each node by the same parameter
$uv=s$. This gives an equivariant smoothing. The cyclic
action remains free after shrinking the parameter disc:
the fixed locus is closed in the proper total family and
has empty central fibre. Thus, for $s\ne0$, the quotient
is an actual etale cyclic cover
$\widetilde C_s\to C_s$ with genera seven and three.
Both central curves are of compact type, so their relative
Jacobians extend as smooth families of polarized abelian varieties and the norm map
extends between them.

Write $J=J(D)$. At the central fibre these maps are the
explicit maps
\[
 J(\widetilde C_0)=J^3\times E'
   \longrightarrow J(C_0)=J\times E,
 \qquad (a,b,c,e)\longmapsto
              (a+b+c,\operatorname{Nm}_{E'/E}e).
 \tag{F1S8.1}
\]
The elliptic norm is an isogeny, so the identity component
of the kernel is
$P_0=\{(a,b,c)\in J^3:a+b+c=0\}$.
The isomorphism
\[
 J^2\xrightarrow{\sim}P_0,\qquad
 (x,y)\longmapsto(x,-y,y-x)
 \tag{F1S8.2}
\]
pulls back the product principal polarization of $J^3$ to
$2\omega_1-\omega_\sigma+2\omega_2$.
Choose the cyclic generator to permute
$(a,b,c)\mapsto(c,a,b)$. Under (F1S8.2) it becomes
$(x,y)\mapsto(y-x,-x)$, exactly the matrix $T$ in
(F1S7.1). Thus this family of actual Pryms specializes to
the same polarized $K$ abelian fourfold from which the
2007 Schoen smoothing is constructed.

\subsection{Status of the original 1988 Prym norm-fibre surface}

Schoen, \emph{Hodge classes on self-products of a variety with an
automorphism}, Compositio Math. 65 (1988), 3--32, Theorem 3.2,
constructs algebraic classes through the cyclic-cover norm-fibre
geometry. The construction reaches its Prym via the Abel map
from effective divisors of degree four on the genus-seven cover.
In the retrieved proof the needed cycle is a proper image; a
specific smooth embedded reduced component together with its
normal bundle has not been identified in this follow-up.
Explicitly, for $\pi:\widetilde C\to C$, the surface source
lies in
$T=(\pi^{(4)})^{-1}|K_C|\subset\operatorname{Sym}^4\widetilde C$.
The Abel map $D\mapsto\mathcal O_{\widetilde C}(D)$ maps it
to the norm fibre over $K_C$. Choosing a point in a component
of that fibre identifies the component with a Prym torsor,
then with $P$ by translation. This is the actual morphism
whose reduced image requires the smoothness and normal-bundle
calculation. The ambient degeneration (F1S8.1)--(F1S8.2)
does not compute the flat limit of this Abel image or an
isomorphism of it with a smooth Schoen member.
The original ATTACK Theorem E is an effective-cycle assertion
and likewise does not specify that embedding.

Van Geemen's introduction explicitly relates Schoen's 2007
construction to a case of the 1988 result. Equations
(F1S5.1)--(F1S5.4) establish its precise $K$ action, polarization,
and Weil class. They do not provide a morphism from the original
norm-fibre component to that smooth member. HD-F1-S08 supplies
an actual comparison of the polarized ambient Prym families.
The computations
HD-F1-S05 and HD-F1-S06 are therefore attached to the stated
smooth Schoen family and its actual multiplication pullbacks.
For the particular original Prym norm-fibre image, Q-T7-FIXED
retains the gate of proving a smooth embedding and calculating
its map (F1S1.1), or its cokernel (F1S1.2) when applicable.
No claimed identification of the two surface constructions is
used in any proved result above.

\paragraph{Resolution of this former gate.}
HD-F1-S09, T07 and S10 below now supply the original norm-image
surface and its conormal map. For every connected etale triple cover
of a nonhyperelliptic genus-three curve the reduced Abel image is
smooth, has degree36, and its requested map has rank12 and zero
cokernel. The preceding status paragraph records the earlier
ambient-only comparison; this later proof closes that gate in the
explicit nonhyperelliptic scope.
